课程教学中的学习记录常常分散在资源、练习与课堂活动中。学生完成一道题后，需要知道困难来自哪里、下一步适合学习什么；教师掌握班级成绩后，还需要理解不同学生的错误，并将判断转为具体的教学安排。学脉围绕课程学习、个性化辅导与教学协作，将课前材料准备、课中问题互动和课后练习反馈组织在同一平台。学生通过知识网络连接讲义、视频、交互实验和习题，教师通过学情证据组织课程、评阅与任务，双方沿同一课程对象持续协作。数据结构与Python课程提供可直接操作的演示内容，新的主题也能够通过课程建设流程接入。

本项目将个性化支持落实为能够继续执行的课程活动，以共享学习证据连接五项方法。知识追踪区分独立作答、提示帮助、相关重复和时间变化，表达知识状态及证据充分程度；路径规划把个人目标、教师任务与到期复习纳入同一预算，通过先修约束安排活动顺序。学生辅导在诊断、教学表征和校验之间传递状态，用课程资料解释问题，再通过独立理解检查调整后续策略。班级教学分层依据错因与记录支持度形成临时关注群组，连接薄弱知识和练习；教师助手综合教学目标与授权班级证据，形成可编辑、可校验的方案。本项目将模型生成、确定性评分和教师确认分别交给适合的执行主体，使解释能够进入学习，教学产物能够保存、发布并追踪结果。平台采用React、FastAPI和SQLite实现，统一维护课程资料、引用与活动记录。

验证覆盖本地方法对照、真实模型任务和师生业务三个层次。1200条学习轨迹用于比较证据修正与消融；180个独立先修图形成720种预算与状态场景，路径方法的先修及预算违规均为0，相对同约束可行贪心的平均效用提高\PathGain\%。48个不同规模班级的184项任务匹配均对应主要薄弱知识，三课程180个评测问题的前五检索来源覆盖全部目标材料。真实DeepSeek的\RealTasks{}个独立任务采用配对提示和相同业务约束，进一步核查理解检查、教学草稿与保存结果。这些结果说明，平台能够在有限学习时间内形成可执行安排，将学情判断落到具体任务，并保留师生核对与调整的依据。学脉由此形成连接学习困难、个性化活动与教学反馈的完整支持过程。
